\section{Injective, Surjective, Bijective}
\chapterline

Once you view a function as a mapping between two sets, three natural questions arise about its structural geometry:
\begin{enumerate}[label=\arabic*.]
\item \textbf{Do distinct inputs always land on distinct outputs?} Or do some paths merge and destroy information along the way?
\item \textbf{Does the mapping reach every corner of the codomain?} Or are some target elements permanently out of reach?
\item \textbf{Can the process be cleanly run in reverse?} Can we trace each output uniquely backward to the input that produced it?
\end{enumerate}

These three questions correspond to the three classical pillars of mapping theory: \emph{injectivity}, \emph{surjectivity}, and \emph{bijectivity}.

\begin{historybox}
Before the late nineteenth century, mathematicians had no universal vocabulary for classifying mappings. Equations were solved by algebraic manipulation, and invertibility was handled on an ad-hoc basis. The language of injective, surjective, and bijective maps emerged when Richard Dedekind and Georg Cantor confronted the profound mystery of the infinite.

In his groundbreaking 1888 treatise, \emph{Was sind und was sollen die Zahlen?} (``What are numbers and what should they be?''), Dedekind sought to define the natural numbers and infinity without relying on physical time, geometric space, or counting. His radical breakthrough was to define infinite sets purely through mapping properties: \emph{a set $S$ is infinite if and only if there exists an injective function $f: S \to S$ that is not surjective.} An infinite set is an object that can be squeezed into a proper subset of itself without any collisions! 

Simultaneously, Cantor employed bijections to accomplish what had previously seemed impossible: measuring the size (cardinality) of infinite collections without counting them one by one. By demonstrating that two sets share the same size if and only if there exists a bijective pairing between them, Cantor transformed functions from computational tools into the fundamental scales used to weigh infinity.
\end{historybox}

\begin{center}
\begin{tikzpicture}[scale=0.88, >=Latex]
  % Panel 1: Injective
  \begin{scope}[xshift=0cm]
    \node[anchor=south, font=\footnotesize\bfseries] at (2.0, 3.4) {1. Injective (No Collisions)};
    \draw[rounded corners=4pt, fill=gray!3, draw=gray!25] (-0.3, -0.6) rectangle (4.3, 3.3);
    
    % Sets
    \draw[thick, sourceblue, fill=sourceblue!12] (0.8, 1.3) ellipse (0.7cm and 1.3cm);
    \draw[thick, targetgreen, fill=targetgreen!12] (3.2, 1.3) ellipse (0.7cm and 1.3cm);
    
    \node[sourceblue!90!black, font=\scriptsize\bfseries] at (0.8, 2.3) {$A$};
    \node[targetgreen!90!black, font=\scriptsize\bfseries] at (3.2, 2.3) {$B$};
    
    % Elements A
    \node[circle, fill=sourceblue, inner sep=1.5pt, label={left:\scriptsize $a_1$}] (i1) at (0.8, 1.7) {};
    \node[circle, fill=sourceblue, inner sep=1.5pt, label={left:\scriptsize $a_2$}] (i2) at (0.8, 0.9) {};
    
    % Elements B
    \node[circle, fill=targetgreen, inner sep=1.5pt, label={right:\scriptsize $b_1$}] (ib1) at (3.2, 1.9) {};
    \node[circle, fill=targetgreen, inner sep=1.5pt, label={right:\scriptsize $b_2$}] (ib2) at (3.2, 1.3) {};
    \node[circle, fill=gray!60, inner sep=1.5pt, label={right:\scriptsize $b_3$}] (ib3) at (3.2, 0.7) {};
    
    % Arrows
    \draw[->, thick, sourceblue!85!black] (i1) -- (ib1);
    \draw[->, thick, sourceblue!85!black] (i2) -- (ib2);
    
    \node[below, font=\scriptsize, align=center] at (2.0, -0.05) {Distinct inputs $\implies$\\distinct outputs};
  \end{scope}

  % Panel 2: Surjective
  \begin{scope}[xshift=5.5cm]
    \node[anchor=south, font=\footnotesize\bfseries] at (2.0, 3.4) {2. Surjective (Full Coverage)};
    \draw[rounded corners=4pt, fill=gray!3, draw=gray!25] (-0.3, -0.6) rectangle (4.3, 3.3);
    
    % Sets
    \draw[thick, sourceblue, fill=sourceblue!12] (0.8, 1.3) ellipse (0.7cm and 1.3cm);
    \draw[thick, targetgreen, fill=targetgreen!12] (3.2, 1.3) ellipse (0.7cm and 1.3cm);
    
    \node[sourceblue!90!black, font=\scriptsize\bfseries] at (0.8, 2.3) {$A$};
    \node[targetgreen!90!black, font=\scriptsize\bfseries] at (3.2, 2.3) {$B$};
    
    % Elements A
    \node[circle, fill=sourceblue, inner sep=1.5pt, label={left:\scriptsize $a_1$}] (s1) at (0.8, 2.0) {};
    \node[circle, fill=sourceblue, inner sep=1.5pt, label={left:\scriptsize $a_2$}] (s2) at (0.8, 1.3) {};
    \node[circle, fill=sourceblue, inner sep=1.5pt, label={left:\scriptsize $a_3$}] (s3) at (0.8, 0.6) {};
    
    % Elements B
    \node[circle, fill=targetgreen, inner sep=1.5pt, label={right:\scriptsize $b_1$}] (sb1) at (3.2, 1.7) {};
    \node[circle, fill=targetgreen, inner sep=1.5pt, label={right:\scriptsize $b_2$}] (sb2) at (3.2, 0.9) {};
    
    % Arrows
    \draw[->, thick, sourceblue!85!black] (s1) -- (sb1);
    \draw[->, thick, sourceblue!85!black] (s2) -- (sb1);
    \draw[->, thick, sourceblue!85!black] (s3) -- (sb2);
    
    \node[below, font=\scriptsize, align=center] at (2.0, -0.05) {Every target hit;\\collisions allowed};
  \end{scope}

  % Panel 3: Bijective
  \begin{scope}[xshift=11.0cm]
    \node[anchor=south, font=\footnotesize\bfseries] at (2.0, 3.4) {3. Bijective (1-to-1 Pairing)};
    \draw[rounded corners=4pt, fill=gray!3, draw=gray!25] (-0.3, -0.6) rectangle (4.3, 3.3);
    
    % Sets
    \draw[thick, sourceblue, fill=sourceblue!12] (0.8, 1.3) ellipse (0.7cm and 1.3cm);
    \draw[thick, quotientorange, fill=quotientorange!12] (3.2, 1.3) ellipse (0.7cm and 1.3cm);
    
    \node[sourceblue!90!black, font=\scriptsize\bfseries] at (0.8, 2.3) {$A$};
    \node[quotientorange!90!black, font=\scriptsize\bfseries] at (3.2, 2.3) {$B$};
    
    % Elements A
    \node[circle, fill=sourceblue, inner sep=1.5pt, label={left:\scriptsize $a_1$}] (b1) at (0.8, 1.9) {};
    \node[circle, fill=sourceblue, inner sep=1.5pt, label={left:\scriptsize $a_2$}] (b2) at (0.8, 1.3) {};
    \node[circle, fill=sourceblue, inner sep=1.5pt, label={left:\scriptsize $a_3$}] (b3) at (0.8, 0.7) {};
    
    % Elements B
    \node[circle, fill=quotientorange, inner sep=1.5pt, label={right:\scriptsize $b_1$}] (bb1) at (3.2, 1.9) {};
    \node[circle, fill=quotientorange, inner sep=1.5pt, label={right:\scriptsize $b_2$}] (bb2) at (3.2, 1.3) {};
    \node[circle, fill=quotientorange, inner sep=1.5pt, label={right:\scriptsize $b_3$}] (bb3) at (3.2, 0.7) {};
    
    % Arrows
    \draw[<->, thick, structurepurple!90!black] (b1) -- (bb1);
    \draw[<->, thick, structurepurple!90!black] (b2) -- (bb2);
    \draw[<->, thick, structurepurple!90!black] (b3) -- (bb3);
    
    \node[below, font=\scriptsize, align=center] at (2.0, -0.05) {Perfect matching;\\two-way invertible};
  \end{scope}
\end{tikzpicture}
\end{center}

\subsection{Injectivity: Collision-Free Mappings}

\begin{definition}[Injective]
A function $f: A \to B$ is \textbf{injective} (or \textbf{one-to-one}) if for all $a_1, a_2 \in A$,
\[
f(a_1) = f(a_2) \implies a_1 = a_2.
\]
Equivalently, by contrapositive:
\[
a_1 \neq a_2 \implies f(a_1) \neq f(a_2).
\]
\end{definition}

An injective map preserves distinctions. If you take two genuinely distinct inputs, their outputs remain distinct. No two input arrows ever crash into the same destination point.

\begin{example}
The linear function $f: \R \to \R$ defined by $f(x) = 3x - 5$ is injective.
To prove this, we follow the definition directly: assume $f(a_1) = f(a_2)$ and deduce $a_1 = a_2$:
\[
3a_1 - 5 = 3a_2 - 5 \implies 3a_1 = 3a_2 \implies a_1 = a_2.
\]
Because $a_1 = a_2$ is inescapable, $f$ never merges inputs.
\end{example}

\begin{example}[A Collision Failure]
The function $g: \R \to \R$ given by $g(x) = x^2$ fails to be injective. A single concrete counterexample suffices: $g(2) = 4$ and $g(-2) = 4$, yet $2 \neq -2$. The sign information is erased.
\end{example}

\subsection{Surjectivity: Covering the Target Space}

\begin{definition}[Surjective]
A function $f: A \to B$ is \textbf{surjective} (or \textbf{onto}) if for every element $b \in B$, there exists at least one element $a \in A$ such that
\[
f(a) = b.
\]
In set language: $\operatorname{im}(f) = B$.
\end{definition}

Surjectivity asks whether the function leaves any empty pockets in the codomain. If you point to any target element $b$ whatsoever, there must be an incoming arrow hitting it.

\begin{example}
The cubic function $f: \R \to \R$ given by $f(x) = x^3$ is surjective.
To prove this, let $y \in \R$ be an arbitrary element of the codomain. We must exhibit an input $x \in \R$ satisfying $f(x) = y$.
We choose $x = \sqrt[3]{y}$. Since every real number has a unique real cube root, $x \in \R$, and
\[
f(x) = (\sqrt[3]{y})^3 = y.
\]
Thus every target $y$ is hit, and $f$ is surjective.
\end{example}

\begin{example}[An Uncovered Target]
The function $g: \R \to \R$ given by $g(x) = x^2$ is not surjective. Consider $y = -1 \in \R$. Since $x^2 \ge 0$ for all real numbers $x$, there is no real input satisfying $x^2 = -1$. The negative portion of the codomain is entirely deserted.
\end{example}

\subsection{Bijectivity: The Reversible Pairing}

\begin{definition}[Bijective]
A function $f: A \to B$ is \textbf{bijective} (or a \textbf{one-to-one correspondence}) if it is both injective and surjective.
\end{definition}

When a function is bijective, every single target in $B$ receives \emph{exactly one} arrow from $A$:
\begin{itemize}
\item Because $f$ is surjective, every target receives \textbf{at least one} arrow.
\item Because $f$ is injective, no target receives \textbf{more than one} arrow.
\end{itemize}
This establishes a perfect lock-and-key pairing between the elements of $A$ and the elements of $B$.

\begin{theorem}
A function $f: A \to B$ is bijective if and only if there exists an inverse function $g: B \to A$ such that
\[
g(f(a)) = a \quad \text{for all } a \in A, \qquad \text{and} \qquad f(g(b)) = b \quad \text{for all } b \in B.
\]
\end{theorem}

We write $g = f^{-1}$. The arrow diagram can be completely reversed: you simply turn every arrow around, and the resulting system is itself a well-defined function from $B$ back to $A$.

\begin{intuitionbox}
For finite sets, injectivity says $|A| \le |B|$ (the Pigeonhole Principle: you cannot fit $n$ birds into fewer than $n$ holes without sharing). Surjectivity says $|A| \ge |B|$ (you need at least as many birds as nests to fill every nest). Therefore, bijectivity implies $|A| = |B|$.

For infinite sets, this exact principle resolves a paradox that stalled mathematics for centuries. In 1638, Galileo Galilei noticed in \emph{Two New Sciences} that every positive integer $n$ can be paired with its square $n^2$ ($1 \leftrightarrow 1, 2 \leftrightarrow 4, 3 \leftrightarrow 9$). This mapping is an exact bijection! Yet the squares become sparser and sparser, so most integers are not squares. Trapped by Euclid's classical axiom that ``the whole is greater than the part,'' Galileo concluded that one cannot speak of infinite sets being greater, lesser, or equal to one another. Dedekind and Cantor reversed this perspective: Euclid's axiom holds only for finite sets; for infinite sets, bijectivity \emph{is} the definition of equal cardinality!
\end{intuitionbox}

\subsection{Standard Proof Blueprints}

When writing proofs, follow these time-tested templates:

\begin{focusbox}
\textbf{To prove that $f: A \to B$ is injective:}
\begin{enumerate}
\item Begin: ``Let $a_1, a_2 \in A$ and assume $f(a_1) = f(a_2)$.''
\item Use algebra and the definition of $f$ to unpack the equation.
\item End: ``Therefore $a_1 = a_2$. Hence $f$ is injective.''
\end{enumerate}

\textbf{To prove that $f: A \to B$ is surjective:}
\begin{enumerate}
\item Begin: ``Let $b \in B$ be an arbitrary element of the codomain.''
\item (On scratch paper, solve $f(a) = b$ for $a$.)
\item In your written proof: Define $a$, verify that $a \in A$, and evaluate: ``$f(a) = \dots = b$.''
\item End: ``Since $b \in B$ was arbitrary, $f$ is surjective.''
\end{enumerate}
\end{focusbox}

\begin{labbox}
Consider $f: \R \setminus \{2\} \to \R \setminus \{3\}$ given by
\[
f(x) = \frac{3x - 1}{x - 2}.
\]
\textbf{1. Injectivity:} Assume $f(x_1) = f(x_2)$:
\[
\frac{3x_1 - 1}{x_1 - 2} = \frac{3x_2 - 1}{x_2 - 2} \implies (3x_1 - 1)(x_2 - 2) = (3x_2 - 1)(x_1 - 2).
\]
Expanding both sides:
\[
3x_1 x_2 - 6x_1 - x_2 + 2 = 3x_1 x_2 - x_1 - 6x_2 + 2.
\]
Subtracting common terms yields $-6x_1 - x_2 = -x_1 - 6x_2$, which simplifies to $5x_2 = 5x_1$, so $x_1 = x_2$.

\textbf{2. Surjectivity:} Given an arbitrary $y \in \R \setminus \{3\}$, solve $y = \frac{3x - 1}{x - 2}$ for $x$:
\[
y(x - 2) = 3x - 1 \implies yx - 2y = 3x - 1 \implies x(y - 3) = 2y - 1 \implies x = \frac{2y - 1}{y - 3}.
\]
Because $y \neq 3$, the denominator is non-zero. If $x = 2$, then $2y - 1 = 2(y - 3) = 2y - 6$, leading to $-1 = -6$, an impossibility. Thus $x \in \R \setminus \{2\}$. Computing $f(x)$ confirms $f(x) = y$.
\end{labbox}

\subsection{Combinatorial Proof Weapons: The Pigeonhole Principle and Double Counting}

When we restrict our attention from arbitrary spaces to finite collections, the mapping properties of injectivity, surjectivity, and bijectivity transform into concrete arithmetic laws.

Consider two non-empty finite sets $A$ and $B$, with sizes $|A| = m$ and $|B| = n$.
An injective function $f: A \to B$ demands that every input element $a \in A$ receives its own distinct destination $f(a) \in B$. If each of the $m$ inputs requires an exclusive target, there must be at least as many targets as inputs:
\[
f \text{ is injective} \implies |A| \le |B|.
\]
Now state the contrapositive of this implication:
\[
|A| > |B| \implies f \text{ is not injective}.
\]
This simple statement is one of the sharpest weapons in discrete mathematics. If you attempt to map a larger set into a strictly smaller set, an output collision is mathematically unavoidable: there must exist two distinct inputs $a_1 \neq a_2 \in A$ that land on the exact same output $f(a_1) = f(a_2)$. 

In 1834, Peter Gustav Lejeune Dirichlet formalized this observation as the \emph{Schubfachprinzip} (``drawer principle''), known in English as the \textbf{Pigeonhole Principle}.

\subsubsection{The Generalized Pigeonhole Principle}

What happens when the domain is not just slightly larger than the codomain, but many times larger? If $10$ items are distributed among $3$ boxes, at least one box must receive not merely $2$ items, but at least $4$.

\begin{theorem}[Generalized Pigeonhole Principle]
Let $A$ and $B$ be non-empty finite sets with $|A| = m$ and $|B| = n$. For any function $f: A \to B$, there exists at least one element $b \in B$ whose preimage contains at least $\lceil m/n \rceil$ elements:
\[
|f^{-1}(\{b\})| \ge \left\lceil \frac{m}{n} \right\rceil.
\]
\end{theorem}

\begin{proof}
We proceed by contradiction. Suppose the conclusion is false: every target element $b \in B$ satisfies $|f^{-1}(\{b\})| \le \lceil m/n \rceil - 1$.
Because every element of $A$ maps to a unique element of $B$, the fibers $\{f^{-1}(\{b\})\}_{b \in B}$ partition the domain $A$. Summing the sizes of these fibers over all $n$ targets in $B$:
\[
m = |A| = \sum_{b \in B} |f^{-1}(\{b\})| \le \sum_{b \in B} \left( \left\lceil \frac{m}{n} \right\rceil - 1 \right) = n \left( \left\lceil \frac{m}{n} \right\rceil - 1 \right).
\]
By definition of the ceiling function, $\lceil x \rceil - 1 < x$ for any real number $x$. Setting $x = m/n$, we obtain:
\[
n \left( \left\lceil \frac{m}{n} \right\rceil - 1 \right) < n \cdot \frac{m}{n} = m.
\]
Combining these inequalities gives $m < m$, an impossibility. Therefore, at least one target $b \in B$ must satisfy $|f^{-1}(\{b\})| \ge \lceil m/n \rceil$.
\end{proof}

\begin{intuitionbox}
Think of the distribution of items as water settling across a collection of vertical tubes. To minimize the maximum height in any single tube, you must spread the fluid as flatly as possible. The flattest possible partition of $m$ items across $n$ boxes assigns either $\lfloor m/n \rfloor$ or $\lceil m/n \rceil$ items to each box. 

For $7$ items placed into $3$ boxes (illustrated in Figure~\ref{13:fig-pigeonhole-doublecount}, Panel~A), the flattest possible spread is $(3, 2, 2)$. If you attempted to keep every box at or below $2$ items, the total capacity of all three boxes combined would be $3 \times 2 = 6$, leaving the 7th item without a place. Even the smoothest possible distribution forces at least one box to reach height $\lceil 7/3 \rceil = 3$.
\end{intuitionbox}

\begin{figure}[htbp]
\centering
\begin{tikzpicture}[scale=0.85, >=Latex]
  % ==========================================
  % Panel A: Generalized Pigeonhole Principle
  % ==========================================
  \begin{scope}[xshift=0cm]
    \node[anchor=south, font=\footnotesize\bfseries] at (3.5, 5.0) {Panel A: Generalized Pigeonhole};
    \node[font=\scriptsize, text=gray!80!black] at (3.5, 4.7) {$7$ items placed into $3$ target boxes};
    \draw[rounded corners=4pt, fill=gray!3, draw=gray!25] (-0.3, -1.25) rectangle (7.4, 4.55);

    % Domain container A
    \draw[thick, sourceblue, fill=sourceblue!10, rounded corners=4pt] (0.0, 0.3) rectangle (1.8, 4.25);
    \node[sourceblue!90!black, font=\scriptsize\bfseries] at (0.9, 4.0) {Domain $A$};
    \node[sourceblue!75!black, font=\tiny] at (0.9, 3.7) {($7$ items)};

    % 7 Domain nodes
    \node[circle, fill=sourceblue, inner sep=1.6pt, label={left:\scriptsize $x_1$}] (x1) at (0.9, 3.30) {};
    \node[circle, fill=sourceblue, inner sep=1.6pt, label={left:\scriptsize $x_2$}] (x2) at (0.9, 2.80) {};
    \node[circle, fill=sourceblue, inner sep=1.6pt, label={left:\scriptsize $x_3$}] (x3) at (0.9, 2.30) {};
    \node[circle, fill=sourceblue, inner sep=1.6pt, label={left:\scriptsize $x_4$}] (x4) at (0.9, 1.80) {};
    \node[circle, fill=sourceblue, inner sep=1.6pt, label={left:\scriptsize $x_5$}] (x5) at (0.9, 1.35) {};
    \node[circle, fill=sourceblue, inner sep=1.6pt, label={left:\scriptsize $x_6$}] (x6) at (0.9, 0.90) {};
    \node[circle, fill=sourceblue, inner sep=1.6pt, label={left:\scriptsize $x_7$}] (x7) at (0.9, 0.50) {};

    % Box 1 (Overloaded, holds 3)
    \draw[thick, quotientorange!85!black, fill=quotientorange!15, rounded corners=3pt] (3.4, 2.9) rectangle (7.1, 4.25);
    \draw[quotientorange!50, dashed] (5.6, 2.9) -- (5.6, 4.25);
    \node[anchor=north west, font=\scriptsize\bfseries, quotientorange!95!black] at (3.55, 4.15) {Box $B_1$};
    \node[anchor=north west, font=\tiny\bfseries, quotientorange!85!black] at (3.55, 3.65) {$\mathbf{3}$ items};
    \node[anchor=north west, font=\tiny, quotientorange!85!black] at (3.55, 3.30) {($\ge \lceil 7/3 \rceil$)};
    % Tokens inside Box 1
    \node[circle, fill=quotientorange!90!black, inner sep=1.6pt] (t1) at (6.35, 3.90) {};
    \node[circle, fill=quotientorange!90!black, inner sep=1.6pt] (t2) at (6.35, 3.55) {};
    \node[circle, fill=quotientorange!90!black, inner sep=1.6pt] (t3) at (6.35, 3.20) {};

    % Box 2 (Holds 2)
    \draw[thick, targetgreen!80!black, fill=targetgreen!12, rounded corners=3pt] (3.4, 1.6) rectangle (7.1, 2.7);
    \draw[targetgreen!50, dashed] (5.6, 1.6) -- (5.6, 2.7);
    \node[anchor=north west, font=\scriptsize\bfseries, targetgreen!90!black] at (3.55, 2.55) {Box $B_2$};
    \node[anchor=north west, font=\tiny, targetgreen!80!black] at (3.55, 2.10) {$2$ items};
    % Tokens inside Box 2
    \node[circle, fill=targetgreen!90!black, inner sep=1.6pt] (t4) at (6.35, 2.35) {};
    \node[circle, fill=targetgreen!90!black, inner sep=1.6pt] (t5) at (6.35, 1.95) {};

    % Box 3 (Holds 2)
    \draw[thick, targetgreen!80!black, fill=targetgreen!12, rounded corners=3pt] (3.4, 0.35) rectangle (7.1, 1.45);
    \draw[targetgreen!50, dashed] (5.6, 0.35) -- (5.6, 1.45);
    \node[anchor=north west, font=\scriptsize\bfseries, targetgreen!90!black] at (3.55, 1.30) {Box $B_3$};
    \node[anchor=north west, font=\tiny, targetgreen!80!black] at (3.55, 0.85) {$2$ items};
    % Tokens inside Box 3
    \node[circle, fill=targetgreen!90!black, inner sep=1.6pt] (t6) at (6.35, 1.1) {};
    \node[circle, fill=targetgreen!90!black, inner sep=1.6pt] (t7) at (6.35, 0.7) {};

    % Mapping arrows
    \draw[->, thick, sourceblue!85!black] (x1) -- (3.4, 3.90);
    \draw[->, thick, sourceblue!85!black] (x2) -- (3.4, 3.55);
    \draw[->, thick, sourceblue!85!black] (x3) -- (3.4, 3.20);
    \draw[->, thick, sourceblue!85!black] (x4) -- (3.4, 2.35);
    \draw[->, thick, sourceblue!85!black] (x5) -- (3.4, 1.95);
    \draw[->, thick, sourceblue!85!black] (x6) -- (3.4, 1.10);
    \draw[->, thick, sourceblue!85!black] (x7) -- (3.4, 0.70);

    \node[below, font=\scriptsize, align=center] at (3.5, -0.05) {Flattest spread: $3 + 2 + 2 = 7$\\[2pt]Capacity if each box $\le 2$: $2 \times 3 = 6 < 7$\\[2pt]\textbf{Forces at least one box to hold $\ge 3$}};
  \end{scope}

  % ==========================================
  % Panel B: Double Counting Incidence Grid
  % ==========================================
  \begin{scope}[xshift=8.6cm]
    \node[anchor=south, font=\footnotesize\bfseries] at (3.3, 5.0) {Panel B: Double Counting Grid};
    \node[font=\scriptsize, text=gray!80!black] at (3.3, 4.7) {Incidence table $M \in \{0,1\}^{4 \times 4}$};
    \draw[rounded corners=4pt, fill=gray!3, draw=gray!25] (-0.6, -1.25) rectangle (7.1, 4.55);

    % Column headers
    \node[targetgreen!90!black, font=\scriptsize\bfseries] at (2.35, 4.25) {Targets $b_j \in B$ (Columns)};
    \node[font=\scriptsize\bfseries] at (1.0, 3.9) {$b_1$};
    \node[font=\scriptsize\bfseries] at (1.9, 3.9) {$b_2$};
    \node[font=\scriptsize\bfseries] at (2.8, 3.9) {$b_3$};
    \node[font=\scriptsize\bfseries] at (3.7, 3.9) {$b_4$};
    \node[targetgreen!90!black, font=\scriptsize\bfseries] at (5.05, 3.9) {Row $\sum$};

    % Row headers
    \node[font=\scriptsize\bfseries] at (-0.1, 3.25) {$a_1$};
    \node[font=\scriptsize\bfseries] at (-0.1, 2.55) {$a_2$};
    \node[font=\scriptsize\bfseries] at (-0.1, 1.85) {$a_3$};
    \node[font=\scriptsize\bfseries] at (-0.1, 1.15) {$a_4$};

    % Grid entries (Row 1): [1, 1, 0, 1] -> 3
    \draw[fill=sourceblue!14, draw=sourceblue!40, rounded corners=2pt] (0.6, 2.95) rectangle (1.4, 3.55);
    \node[sourceblue!90!black, font=\small\bfseries] at (1.0, 3.25) {1};
    \draw[fill=sourceblue!14, draw=sourceblue!40, rounded corners=2pt] (1.5, 2.95) rectangle (2.3, 3.55);
    \node[sourceblue!90!black, font=\small\bfseries] at (1.9, 3.25) {1};
    \draw[fill=gray!6, draw=gray!25, rounded corners=2pt] (2.4, 2.95) rectangle (3.2, 3.55);
    \node[gray!50!black, font=\small] at (2.8, 3.25) {0};
    \draw[fill=sourceblue!14, draw=sourceblue!40, rounded corners=2pt] (3.3, 2.95) rectangle (4.1, 3.55);
    \node[sourceblue!90!black, font=\small\bfseries] at (3.7, 3.25) {1};
    \draw[fill=targetgreen!14, draw=targetgreen!45, rounded corners=2pt] (4.5, 2.95) rectangle (5.6, 3.55);
    \node[targetgreen!90!black, font=\small\bfseries] at (5.05, 3.25) {3};

    % Grid entries (Row 2): [0, 1, 1, 0] -> 2
    \draw[fill=gray!6, draw=gray!25, rounded corners=2pt] (0.6, 2.25) rectangle (1.4, 2.85);
    \node[gray!50!black, font=\small] at (1.0, 2.55) {0};
    \draw[fill=sourceblue!14, draw=sourceblue!40, rounded corners=2pt] (1.5, 2.25) rectangle (2.3, 2.85);
    \node[sourceblue!90!black, font=\small\bfseries] at (1.9, 2.55) {1};
    \draw[fill=sourceblue!14, draw=sourceblue!40, rounded corners=2pt] (2.4, 2.25) rectangle (3.2, 2.85);
    \node[sourceblue!90!black, font=\small\bfseries] at (2.8, 2.55) {1};
    \draw[fill=gray!6, draw=gray!25, rounded corners=2pt] (3.3, 2.25) rectangle (4.1, 2.85);
    \node[gray!50!black, font=\small] at (3.7, 2.55) {0};
    \draw[fill=targetgreen!14, draw=targetgreen!45, rounded corners=2pt] (4.5, 2.25) rectangle (5.6, 2.85);
    \node[targetgreen!90!black, font=\small\bfseries] at (5.05, 2.55) {2};

    % Grid entries (Row 3): [1, 0, 1, 1] -> 3
    \draw[fill=sourceblue!14, draw=sourceblue!40, rounded corners=2pt] (0.6, 1.55) rectangle (1.4, 2.15);
    \node[sourceblue!90!black, font=\small\bfseries] at (1.0, 1.85) {1};
    \draw[fill=gray!6, draw=gray!25, rounded corners=2pt] (1.5, 1.55) rectangle (2.3, 2.15);
    \node[gray!50!black, font=\small] at (1.9, 1.85) {0};
    \draw[fill=sourceblue!14, draw=sourceblue!40, rounded corners=2pt] (2.4, 1.55) rectangle (3.2, 2.15);
    \node[sourceblue!90!black, font=\small\bfseries] at (2.8, 1.85) {1};
    \draw[fill=sourceblue!14, draw=sourceblue!40, rounded corners=2pt] (3.3, 1.55) rectangle (4.1, 2.15);
    \node[sourceblue!90!black, font=\small\bfseries] at (3.7, 1.85) {1};
    \draw[fill=targetgreen!14, draw=targetgreen!45, rounded corners=2pt] (4.5, 1.55) rectangle (5.6, 2.15);
    \node[targetgreen!90!black, font=\small\bfseries] at (5.05, 1.85) {3};

    % Grid entries (Row 4): [0, 1, 0, 1] -> 2
    \draw[fill=gray!6, draw=gray!25, rounded corners=2pt] (0.6, 0.85) rectangle (1.4, 1.45);
    \node[gray!50!black, font=\small] at (1.0, 1.15) {0};
    \draw[fill=sourceblue!14, draw=sourceblue!40, rounded corners=2pt] (1.5, 0.85) rectangle (2.3, 1.45);
    \node[sourceblue!90!black, font=\small\bfseries] at (1.9, 1.15) {1};
    \draw[fill=gray!6, draw=gray!25, rounded corners=2pt] (2.4, 0.85) rectangle (3.2, 1.45);
    \node[gray!50!black, font=\small] at (2.8, 1.15) {0};
    \draw[fill=sourceblue!14, draw=sourceblue!40, rounded corners=2pt] (3.3, 0.85) rectangle (4.1, 1.45);
    \node[sourceblue!90!black, font=\small\bfseries] at (3.7, 1.15) {1};
    \draw[fill=targetgreen!14, draw=targetgreen!45, rounded corners=2pt] (4.5, 0.85) rectangle (5.6, 1.45);
    \node[targetgreen!90!black, font=\small\bfseries] at (5.05, 1.15) {2};

    % Col sums (Row 5): [2, 3, 2, 3] and Grand Total
    \draw[fill=targetgreen!14, draw=targetgreen!45, rounded corners=2pt] (0.6, 0.15) rectangle (1.4, 0.75);
    \node[targetgreen!90!black, font=\small\bfseries] at (1.0, 0.45) {2};
    \draw[fill=targetgreen!14, draw=targetgreen!45, rounded corners=2pt] (1.5, 0.15) rectangle (2.3, 0.75);
    \node[targetgreen!90!black, font=\small\bfseries] at (1.9, 0.45) {3};
    \draw[fill=targetgreen!14, draw=targetgreen!45, rounded corners=2pt] (2.4, 0.15) rectangle (3.2, 0.75);
    \node[targetgreen!90!black, font=\small\bfseries] at (2.8, 0.45) {2};
    \draw[fill=targetgreen!14, draw=targetgreen!45, rounded corners=2pt] (3.3, 0.15) rectangle (4.1, 0.75);
    \node[targetgreen!90!black, font=\small\bfseries] at (3.7, 0.45) {3};

    % Col sum label on left
    \node[targetgreen!90!black, font=\scriptsize\bfseries] at (-0.1, 0.45) {$\sum$};

    % Grand total cell
    \draw[fill=quotientorange!20, draw=quotientorange!80!black, line width=1.1pt, rounded corners=2pt] (4.5, 0.15) rectangle (5.6, 0.75);
    \node[quotientorange!95!black, font=\small\bfseries] at (5.05, 0.45) {10};

    \node[below, font=\scriptsize, align=center] at (3.3, -0.05) {Row sums: $3 + 2 + 3 + 2 = \mathbf{10}$\\[2pt]Column sums: $2 + 3 + 2 + 3 = \mathbf{10}$\\[2pt]\textbf{Grand Total: } $\sum\text{rows} = \sum\text{cols} = \mathbf{10}$};
  \end{scope}
\end{tikzpicture}
\caption{The two workhorses of finite mapping theory. \textbf{Left}: The Generalized Pigeonhole Principle with $7$ items placed into $3$ boxes. Even the most uniform partition ($3, 2, 2$) forces at least one box to contain $\ge \lceil 7/3 \rceil = 3$ items. \textbf{Right}: Double counting on an incidence matrix $M \in \{0,1\}^{4 \times 4}$. Summing horizontally along rows or vertically down columns yields the exact same total: $\sum_{a} \text{row}(a) = \sum_{b} \text{col}(b) = 10$.}
\label{13:fig-pigeonhole-doublecount}
\end{figure}

The true craft of using the Pigeonhole Principle lies in choosing the pigeonholes. In most non-trivial applications, the boxes are not presented in advance; you must construct them from the underlying mathematical structure.

\begin{example}[Differences Divisible by $n$]
Let $n \in \N$ and let $S \subset \Z$ be any collection of $n+1$ distinct integers. Prove that there exist two distinct integers $x, y \in S$ whose difference $x - y$ is a multiple of $n$.

\textbf{Proof.}
By the division algorithm, dividing any integer by $n$ yields a remainder in the set:
\[
R = \{0, 1, 2, \dots, n-1\}.
\]
The set $R$ contains exactly $n$ possible remainders. These $n$ remainders are our pigeonholes.

Define the mapping $f: S \to R$ by $f(z) = z \bmod n$. The domain $S$ contains $n+1$ elements (the pigeons), while the codomain $R$ contains $n$ elements (the pigeonholes). 

Because $|S| = n + 1 > n = |R|$, the function $f$ cannot be injective. There exist distinct elements $x, y \in S$ such that $f(x) = f(y)$. Hence $x$ and $y$ share the exact same remainder $r \in R$:
\[
x = q_1 n + r \quad \text{and} \quad y = q_2 n + r \quad \text{for some } q_1, q_2 \in \Z.
\]
Subtracting the two expressions:
\[
x - y = (q_1 n + r) - (q_2 n + r) = (q_1 - q_2)n.
\]
Since $q_1 - q_2$ is an integer, $x - y$ is an exact multiple of $n$. 

To see this in action, pick $n = 5$ and choose $6$ arbitrary integers: $S = \{14, 23, 37, 51, 69, 93\}$. Computing each remainder modulo $5$:
\[
14 \equiv 4, \quad 23 \equiv 3, \quad 37 \equiv 2, \quad 51 \equiv 1, \quad 69 \equiv 4, \quad 93 \equiv 3 \pmod 5.
\]
Both $14$ and $69$ land in remainder box $4$; their difference $69 - 14 = 55 = 5 \times 11$ is a multiple of $5$. Both $23$ and $93$ land in remainder box $3$; their difference $93 - 23 = 70 = 5 \times 14$ is also a multiple of $5$. The collision was guaranteed before inspecting a single digit.
\end{example}

\begin{example}[Points in a $2 \times 2$ Square]
Suppose $5$ points are placed inside a $2 \times 2$ square $Q = [0, 2] \times [0, 2]$ in the plane. Prove that at least two of the points must lie within a Euclidean distance of at most $\sqrt{2}$ of each other.

\textbf{Proof.}
We partition the continuous $2 \times 2$ square into four $1 \times 1$ sub-squares:
\[
Q_1 = [0, 1] \times [0, 1], \quad Q_2 = [1, 2] \times [0, 1], \quad Q_3 = [0, 1] \times [1, 2], \quad Q_4 = [1, 2] \times [1, 2].
\]
If a point lies on an internal boundary shared by multiple sub-squares, assign it to any one of the adjacent squares arbitrarily. 

We have $5$ points and $4$ sub-squares. By the Pigeonhole Principle, at least one sub-square $Q_k$ must contain at least $\lceil 5/4 \rceil = 2$ points. Let these two points be $P_1 = (x_1, y_1)$ and $P_2 = (x_2, y_2)$.

Because both points lie inside the same $1 \times 1$ square $Q_k$, their horizontal and vertical coordinate separations are bounded:
\[
|x_1 - x_2| \le 1 \qquad \text{and} \qquad |y_1 - y_2| \le 1.
\]
The Euclidean distance $d(P_1, P_2)$ between them is bounded by the diagonal of the unit square:
\[
d(P_1, P_2) = \sqrt{(x_1 - x_2)^2 + (y_1 - y_2)^2} \le \sqrt{1^2 + 1^2} = \sqrt{2} \approx 1.414.
\]
Observe the geometric constraint at work: even if an adversary places the first four points at the extreme outer corners $(0,0), (2,0), (0,2), (2,2)$---maximizing their separations to $2$ and $2\sqrt{2}$---the fifth point cannot escape. Wherever it is placed, it enters a quadrant already containing a corner point, forcing a separation of at most $\sqrt{2}$.
\end{example}

\subsubsection{Double Counting: Counting a 0--1 Grid by Rows vs. by Columns}

A bijection proves that $|A| = |B|$ by building a one-to-one correspondence between two sets. But what if we wish to establish an equality between two complicated combinatorial formulas where an explicit element-by-element matching is elusive?

The technique of \textbf{Double Counting} (or two-way counting) provides the answer. Instead of matching elements directly, we count the size of a single set of relational pairs in two distinct ways.

Let $A$ and $B$ be finite sets, and let $R \subseteq A \times B$ be a relation between them. We encode $R$ as an \emph{incidence matrix} $M \in \{0, 1\}^{|A| \times |B|}$, where each entry indicates whether a pair belongs to $R$:
\[
M_{a, b} = \begin{cases} 1 & \text{if } (a, b) \in R, \\ 0 & \text{if } (a, b) \notin R. \end{cases}
\]
Now ask a simple question: \emph{How many total $1$s are present in the table?}

We can tally the total sum $S$ in two legitimate ways:
\begin{enumerate}
\item \textbf{Row by Row (Horizontal Sum):} Sum the entries across each row first, obtaining the row total $\text{row}(a) = \sum_{b \in B} M_{a, b}$, and then add all row totals together:
\[
S = \sum_{a \in A} \text{row}(a) = \sum_{a \in A} \left( \sum_{b \in B} M_{a, b} \right).
\]
\item \textbf{Column by Column (Vertical Sum):} Sum the entries down each column first, obtaining the column total $\text{col}(b) = \sum_{a \in A} M_{a, b}$, and then add all column totals together:
\[
S = \sum_{b \in B} \text{col}(b) = \sum_{b \in B} \left( \sum_{a \in A} M_{a, b} \right).
\]
\end{enumerate}
Because finite addition is commutative and associative, changing the order of summation leaves the grand total unchanged. We obtain the universal identity:
\[
\sum_{a \in A} \text{row}(a) = \sum_{b \in B} \text{col}(b).
\]
This identity is illustrated in Figure~\ref{13:fig-pigeonhole-doublecount}, Panel~B: summing the rows yields $3 + 2 + 3 + 2 = 10$, and summing the columns yields $2 + 3 + 2 + 3 = 10$. Both roads arrive at the same destination.

\begin{example}[Euler's Handshake Lemma]
In any finite graph $G = (V, E)$, the sum of the degrees of all vertices equals twice the number of edges:
\[
\sum_{v \in V} \deg(v) = 2|E|.
\]

\textbf{Proof.}
Construct an incidence table where rows correspond to vertices $v \in V$ and columns correspond to edges $e \in E$. Set the cell entry to $1$ if vertex $v$ is an endpoint of edge $e$, and $0$ otherwise:
\[
M_{v, e} = \begin{cases} 1 & \text{if } v \in e, \\ 0 & \text{if } v \notin e. \end{cases}
\]
We evaluate the total number of $1$s in $M$ in two ways:
\begin{itemize}
\item \textbf{Summing across rows:} For a fixed vertex $v$, the row sum counts how many edges have $v$ as an endpoint. By definition, this is the degree $\deg(v)$. Summing across all vertices gives $\sum_{v \in V} \deg(v)$.
\item \textbf{Summing down columns:} For a fixed edge $e = \{u, w\}$, the edge connects exactly $2$ endpoints. Thus every single column contains exactly two $1$s. Summing across all $|E|$ edges gives $\sum_{e \in E} 2 = 2|E|$.
\end{itemize}
Equating both expressions yields $\sum_{v \in V} \deg(v) = 2|E|$. 

An immediate structural consequence follows: because $2|E|$ is an even integer, the sum of all vertex degrees must be even. Hence, the number of vertices with odd degree must be even (at any gathering, the number of people who have shaken hands an odd number of times is always even).
\end{example}

\begin{example}[Counting Committee-Chair Pairs]
Let $n, k \in \N$ with $1 \le k \le n$. Prove combinatorially that:
\[
k \binom{n}{k} = n \binom{n-1}{k-1}.
\]

\textbf{Proof.}
Rather than expanding into factorials $\frac{n!}{(k-1)!(n-k)!}$ and canceling terms algebraically, we count a single set in two ways.

Suppose from a pool of $n$ people, we want to form a committee of $k$ members, and from within that committee, select one person to serve as chair. Define the set of valid configurations:
\[
\mathcal{S} = \{(C, p) : C \subseteq P, \, |C| = k, \, p \in C\},
\]
where $P$ is the population of $n$ people, $C$ is a committee of size $k$, and $p \in C$ is the designated chair. We compute $|\mathcal{S}|$ from two viewpoints:
\begin{itemize}
\item \textbf{Viewpoint 1 (Choose committee first, then select chair):}
There are $\binom{n}{k}$ ways to choose the committee $C$ of size $k$ from the $n$ candidates. Once $C$ is fixed, there are $k$ choices for which committee member serves as chair $p$. The number of pairs is:
\[
|\mathcal{S}| = \binom{n}{k} \times k = k \binom{n}{k}.
\]
\item \textbf{Viewpoint 2 (Choose chair first, then fill remaining seats):}
There are $n$ ways to select the chair $p$ from the entire pool of $n$ people. Once the chair is established, the remaining $k-1$ committee members must be selected from the remaining $n-1$ available candidates. This can be done in $\binom{n-1}{k-1}$ ways. The number of pairs is:
\[
|\mathcal{S}| = n \times \binom{n-1}{k-1} = n \binom{n-1}{k-1}.
\]
\end{itemize}
Both viewpoints count the exact same set $\mathcal{S}$ of chair-committee pairs. Therefore, the two expressions must be equal:
\[
k \binom{n}{k} = n \binom{n-1}{k-1}.
\]
Double counting transforms an algebraic manipulation into a structural story about perspectives on a single set.
\end{example}

\subsection{What Mapping Theory Unlocked, Its Honest Limits, and Where We Go Next}

Classifying functions into injections, surjections, and bijections provides the rigorous scaffolding for all of higher mathematics. It gives us:
\begin{itemize}
\item \textbf{Invertibility without guessing:} A function has a two-way inverse if and only if it is a bijection. You no longer have to guess whether an equation can be reversed; checking injectivity and surjectivity decides the matter definitively.
\item \textbf{Isomorphisms:} In abstract algebra (Chapter~21), two groups or vector spaces that look completely different on the surface are recognized as fundamentally identical if there exists a structure-preserving bijection between them.
\item \textbf{Measuring the infinite:} In Chapter~19, we will use bijections to discover the astounding fact that while the integers and rational numbers share the same cardinality ($\aleph_0$), the real numbers form a strictly larger infinity that no bijection can reach.
\end{itemize}

Yet we must acknowledge the honest limitation of this framework: injectivity and surjectivity are strictly \textbf{pointwise} concepts. They inspect single inputs $a$ and single outputs $b$ in isolation. 

In geometry, calculus, and topology, however, spaces are rarely understood one lonely point at a time. We need to know how a mapping stretches, tears, folds, or preserves whole geometric regions, open disks, and continuous curves. If you try to analyze a continuous wave or a gravitational field by tracking every individual point, the notation collapses into chaos.

To understand how functions act on whole collections of points simultaneously, we need to lift our machinery from elements to subsets. In Chapter~14, we study \textbf{Images and Inverse Images}, discovering that pulling backward across a function (preimages) preserves the algebraic fabric of sets in ways that pushing forward can never match.

\subsection{Exercises}

\begin{enumerate}
\item For each function, determine whether it is injective, surjective, or bijective. Provide rigorous proofs or explicit counterexamples:
\begin{enumerate}[label=(\alph*)]
\item $f: \Z \to \Z$ defined by $f(n) = 2n + 3$.
\item $g: \Q \to \Q$ defined by $g(x) = 2x + 3$.
\item $h: \R \to \R$ defined by $h(x) = x^3 - x$.
\end{enumerate}
\item Let $f: A \to B$ and $g: B \to C$ be functions. Prove that if both $f$ and $g$ are injective, their composition $g \circ f: A \to C$ is injective.
\item Prove that if $g \circ f$ is surjective, then $g$ must be surjective. Must $f$ also be surjective? Provide a proof or counterexample.
\item Suppose $n+1$ distinct integers are chosen from the set $\{1, 2, \dots, 2n\}$. Prove that there exist two chosen integers such that one divides the other. (\emph{Hint:} Write each chosen integer in the form $2^k \cdot m$, where $m$ is an odd integer, and identify the pigeonholes.)
\item Use double counting to prove the identity $\sum_{k=1}^n k = \binom{n+1}{2}$ by counting pairs of integers $\{x, y\} \subseteq \{1, 2, \dots, n+1\}$ with $x < y$, grouped by the value of the larger element $y$.
\end{enumerate}

\newpage
